\chapter{Experimentální část}
První experiment byl 
proveden dne 31. října 2025 na~duralové
desce s~rozměry 430~x~513~x~2,75 mm. Byl prováděn 
tzv.~pentest - od snímače AE značky 
Olympus byla lámána tuha
o tloušťce 0,3 mm a tvrdosti 2H (dle~normy ČSN~14584\cite{čsn14584})
ve~vzdálenostech 5, 10, 15, 20, 25, 30, 35 a 40 cm.
Podle normy ČSN~14584\cite{čsn14584} byl snímač 
zatížen silou o velikosti 10 N – bylo použito závaží
o~hmotnosti 1000,03 g.
Data ze snímače AE zpracovávala dvoukanálová měřicí ústředna Dakel. Vlastní
naměřené průběhy byly zobrazovány na notebooku, kde byl nainstalován 
program ZEDO \cite{zedo} pro komunikaci s ústřednou.
\begin{figure}[htbp]
    \centering
    \includegraphics[angle= 180, width=0.75\linewidth]{obrazky/pracoviste_dakel.jpg}
    \caption{Měřicí pracoviště}
    \label{fig:dakel_pracoviste}
\end{figure}
Pro každou zmíněnou vzdálenost od snímače bylo provedeno několik 
tzv.~pentestů. Po dokončení každé sady dat (tj. dat určených pro~jednu 
ze~vzdáleností) byl proveden export. Program ZEDO každý průběh exportoval tak,
že vygeneroval binární soubor se surovými hodnotami AD převodu (*.bin),
textový soubor s informacemi o získaných vzorcích a gnuplot soubor (*.plt) pro~snadné
zobrazení grafů mezi~různými operačními systémy. Vypočtené hodnoty \textbf{RMS}, 
\textbf{ASL}, a \textbf{výkon signálu} byly vloženy do textového souboru, jehož názvu bylo
součástí označení "ae-parameters". Pro tyto parametry platí dle dokumentace \cite{zedo}
následující vztahy:

\begin{equation}
    \label{eq:RMS}
    U_{RMS} = \sqrt{\frac{1}{N}\cdot \sum_{i=1}^{N} u_i^2}
\end{equation}
kde:\\
$U_{RMS}\dots$ napětí efektivní hodnoty signálu (z angl. root mean square) [V]\\
N\dots počet diskrétních vzorků napětí (parametr každého z průběhů) [-]\\
$u_i\dots$ hodnota napětí pro daný vzorek i (spočtená z periferní hodnoty AD převodníku) [V]
\begin{equation}
    \label{eq:ASL}
    U_{ASL} = \frac{1}{N}\cdot \sum_{i=1}^{N} |u_i|
\end{equation}
kde:\\
$U_{ASL}\dots $napětí střední hodnoty signálu (z angl. average signal level) [V]\\

\begin{equation}
    \label{eq:pow}
    P = \frac{1}{N}\cdot \sum_{i=1}^{N} u_i^2
\end{equation}
Ve vztazích \ref{eq:RMS}, \ref{eq:ASL} a \ref{eq:pow} figurují
hodnoty napětí. Tyto hodnoty napětí se dopočítávají ze surové hodnoty
AD převodu dle vztahu (je kromě toho uváděn v textovém souboru každého průběhu):
\begin{equation}
    \label{eq:urceni_ui_z_adc}
    u_i = 10 \cdot \frac{adc}{2^{16}}\cdot 10^{-1 \cdot \frac{Gain}{20}}\hspace{10pt}[V]
\end{equation}
kde:\\
adc\dots periferní hodnota AD převodníku získaná z binárního souboru [-]\\
Gain \dots vlastní zesílení měřicí ústředny (udáno 34 dB) [dB]\\
$2^{16} \dots$ počet kvantizačních úrovní 16bitového AD převodníku (odečtení jedničky zanedbáno) [-]\\

Výpočet parametrů RMS, ASL a výkonu signálu byl proveden pomocí 
skriptu v Pythonu. Z textového souboru byly vyčteny veškeré informace 
o~ukládání surových dat – ústředna Dakel ukládá hodnoty AD převodu jako 
2bytový signed integer s endianitou "little endian" – 
nejprve se do paměti uloží least significant byte (LSB), MSB jako druhý.